Ou seja: demonstrar, através do método básico de análise dimensional, que a complacência ($C$) de uma viga de seção quadrada ($d \times d$) está relacionada ao seu comprimento ($l$) e módulo de elasticidade ($E$).

\textbf{{Resposta:}}

% Identificação de Variáveis e Dimensões
%\subsection*{
Dimensões físicas das variáveis  ({L}: Comprimento; {F}: Força):
%}
\begin{itemize}
    \item {Compliância ($C$):} Deflexão (L) por unidade de carga (F): $[L/F]$.
    \item {Módulo de Elasticidade ($E$):} Força por unidade de área: $[F/L^2]$.
    %Medidas lineares.
    \item {Comprimento ($l$) e Espessura ($d$):}  $[L]$.
\end{itemize}

%Procedimento

%Eliminação da dimensão de Força ($F$). Para que a equação seja racional, as dimensões devem se cancelar. Multiplicamos $C$ por $E$ para eliminar a dimensão $F$:
\begin{equation*}
[C \times E] = \left( \frac{L}{F} \right) \times \left( \frac{F}{L^2} \right) = \frac{1}{L}
\end{equation*}

%Construção do Grupo Adimensional ($\Pi_1$)}. O resultado anterior ($1/L$) precisa ser multiplicado por uma dimensão de comprimento para se tornar um número puro (adimensional). Utilizei a espessura $d$:
\begin{equation*}
\Pi_1 = (CE) \times d = \left( \frac{1}{L} \right) \times L = 1 \text{ (número puro é adimensional)}
\end{equation*}

%Relação de Forma ($\Pi_2$). As variáveis restantes de comprimento ($l$ e $d$) são agrupadas em uma razão para manter a falta de dimensões:
\begin{equation*}
\Pi_2 = \frac{l}{d}
\end{equation*}

%Conclusão
%De acordo com o princípio da homogeneidade dimensional, a função que descreve o fenômeno deve relacionar esses grupos adimensionais:
Portanto: 
%\begin{equation*}
$CEd = f_{CEd}\left(\frac{l}{d} \right)$
%\end{equation*}
, pois o modelo mostra que a flexibilidade da viga é governada pela sua {proporção geométrica} ($l/d$) e pelas {propriedades do material} ($E$), independentemente da escala absoluta do objeto.
\if 0
\section*{Comentários Finais}
\begin{itemize}
    \item \textbf{Por que não considerei $d^2$ em vez de $d$:} Embora a área seja $d^2$, a análise dimensional foca no cancelamento das unidades líquidas. Como o produto $CE$ resultou em $1/L$, basta multiplicar por $d^1$ para "zerar" a dimensão.
    \item \textbf{Significado:} O modelo mostra que a flexibilidade da viga é governada pela sua \textbf{proporção geométrica} ($l/d$) e pelas \textbf{propriedades do material} ($E$), independentemente da escala absoluta do objeto.
\end{itemize}
\fi